\section{Fixed-cut incomplete Gamma tail: differential transcendence, zero density, and the Euler-motive barrier}

\subsection{Status and scope}

This section combines the all-order algebraic-point theorem of companion
paper \textbf{[001]} with the universal inverse-trace identities developed
above.  Accordingly, every arithmetic-independence or transcendence statement
in this section has status \RPASS{}.  Analytic identities are unconditional.
No claim about the individual irrationality or transcendence of \(\gamma\)
or \(\delta\), and no comparison-injectivity claim, is made or used.

Put

\[
 K:=\overline{\mathbf Q}\bigl(e^\alpha:\alpha\in\overline{\mathbf Q}\bigr),
 \qquad E:=E_1(1)=\frac{\delta}{e},
 \qquad \eta:=\Ein(1)=\gamma+E .
\]

Write the Laurent expansion and the upper-tail moments as

\[
 \Gamma(s)=\frac1s+\sum_{k\ge0}g_ks^k,
 \qquad
 T_k:=\frac1{k!}\int_1^\infty e^{-u}\frac{(\log u)^k}{u}\,du .
\]

The split at one proved in {[}001{]} is

\[
 c_k:=g_k-T_k=(-1)^{k+1}\Ein_{k+1}(1),                         \tag{GT-1}
\]

and the family \((c_k)_{k\ge0}\) is jointly algebraically independent over \(K\).

\begin{center}\rule{0.5\linewidth}{0.5pt}\end{center}

\subsection{1. Exact finite gamma-jet field and the tail transcendence-degree bound}

\subsubsection{Lemma 1 (exact gamma-jet field)}

For \(N=0\),

\[
 K(g_0)=K(\gamma).
\]

For every \(N\ge1\),

\[
 \boxed{
 K(g_0,\ldots,g_N)
 =K\!\left(\gamma,\pi^2,\zeta(3),\zeta(5),\ldots,
 \zeta\!\left(2\left\lfloor\frac N2\right\rfloor+1\right)\right).}       \tag{GT-2}
\]

Consequently,

\[
 \operatorname{trdeg}_K K(g_0,\ldots,g_N)
 \le 2+\left\lfloor\frac N2\right\rfloor \qquad(N\ge1).                 \tag{GT-3}
\]

\paragraph{Proof}

The standard identity

\[
 \Gamma(1+s)=
 \exp\!\left(-\gamma s+\sum_{j\ge2}\frac{(-1)^j\zeta(j)}j s^j\right)
\]

shows triangularly that \(g_0,\ldots,g_N\) and
\(\gamma,\zeta(2),\ldots,\zeta(N+1)\) generate the same field: taking the
formal logarithm gives the inverse triangular transformation. The even zeta
values all lie in \(\mathbf Q(\pi^2)\), while the odd indices in
\([3,N+1]\) are precisely

\[
 3,5,\ldots,2\left\lfloor\frac N2\right\rfloor+1,
\]

of which there are \(\lfloor N/2\rfloor\). This proves (2) and (3). Notice
that (3) is only an upper bound: no algebraic independence of \(\gamma\),
\(\pi^2\), or odd zeta values over \(K\) is being assumed. \(\square\)

\subsubsection{Theorem 2 (half-rank lower bound for the fixed-cut tail)}

For every \(N\ge1\),

\[
 \boxed{
 \operatorname{trdeg}_K K(T_0,\ldots,T_N)
 \ge \left\lceil\frac N2\right\rceil-1.}                                  \tag{GT-4}
\]

For \(N=0\), the corresponding lower bound is \(0\). In particular,

\[
 \operatorname{trdeg}_K K(T_0,T_1,\ldots)=\infty.                          \tag{GT-5}
\]

\paragraph{Proof}

By (1),

\[
 K(c_0,\ldots,c_N)\subset
 K(g_0,\ldots,g_N,T_0,\ldots,T_N).
\]

The left side has transcendence degree \(N+1\) over \(K\). Subadditivity
and (3) therefore give

\[
\begin{aligned}
N+1
&\le \operatorname{trdeg}_K K(g_0,\ldots,g_N,T_0,\ldots,T_N)\\
&\le 2+\left\lfloor\frac N2\right\rfloor
 +\operatorname{trdeg}_K K(T_0,\ldots,T_N),
\end{aligned}
\]

which is (4). \(\square\)

The bound is collective. It implies that infinitely many tail moments are
transcendental over \(K\), with lower density at least \(1/2\) among the
indices in the weak counting sense supplied by (4), but it does \textbf{not} name
\(T_0=\delta/e\) as one of them.

\begin{center}\rule{0.5\linewidth}{0.5pt}\end{center}

\subsection{2. Mahler's coefficient-field theorem and a fixed-cut DT theorem}

\subsubsection{The exact Mahler input}

The needed result is the formal-power-series theorem in K. Mahler,
``A remark on algebraic differential equations,'' \emph{Rend. Accad. Naz. Lincei},
Ser. VIII \textbf{50} (1971), 402--412, especially §§3--5
(\href{https://carmamaths.org/resources/mahler/docs/176.pdf}{scan}).

\begin{quote}
\textbf{Mahler's theorem (form used here).} Let \(L\) be any field of
characteristic zero and let \(f(z)=\sum_{n\ge0}f_nz^n\in L[[z]]\). If
\(f\) satisfies a nonzero algebraic differential equation
\[
F(z,f,f',\ldots,f^{(m)})=0,
\qquad F\in L[z,Y_0,\ldots,Y_m],
\]
then the coefficient field \(\mathbf Q(f_0,f_1,\ldots)\) is a finitely
generated field extension of \(\mathbf Q\); in particular, it has finite
transcendence degree over \(\mathbf Q\).
\end{quote}

Rational functions in \(z\) are allowed after clearing denominators. The
statement is formal and includes singular formal solutions. Mahler first
constructs a defining equation over the coefficient field. The
Hurwitz--Popken recurrence in §4 has the eventual form

\[
 A(n)f_{n+1}=B_n(f_0,\ldots,f_n),                                           \tag{GT-6}
\]

where \(A\) is a nonzero polynomial and the coefficients used in all
\(B_n\)'s belong to the finite set of coefficients of the defining equation.
Since a nonzero polynomial has only finitely many integer zeros, all
sufficiently late coefficients are rational functions of finitely many early
coefficients and finitely many equation coefficients. Section 5 records the
result as

\[
 \mathbf Q(f_0,f_1,\ldots)=\mathbf Q(s_1,\ldots,s_r,t),
\]

with \(s_1,\ldots,s_r\) algebraically independent and \(t\) algebraic over
\(\mathbf Q(s_1,\ldots,s_r)\).

\subsubsection{Theorem 3 (fixed-cut upper incomplete Gamma is differentially transcendental)}

Let

\[
 U(s):=\Gamma(s,1)=\int_1^\infty e^{-u}u^{s-1}\,du .
\]

Then

\[
 \boxed{U(s)\text{ is differentially transcendental over }\mathbf C(s).}    \tag{GT-7}
\]

Equivalently, no nonzero polynomial in \(s,U,U',\ldots,U^{(m)}\), with
coefficients in \(\mathbf C\), vanishes identically.

\paragraph{Proof}

Uniform domination on compact subsets of the \(s\)-plane permits
differentiation under the integral, so \(U\) is entire and

\[
 U(s)=\sum_{k\ge0}T_ks^k.                                                   \tag{GT-8}
\]

If \(U\) were differentially algebraic over \(\mathbf C(s)\), clearing
denominators and applying Mahler's theorem with \(L=\mathbf C\) would make
\(\mathbf Q(T_0,T_1,\ldots)\) finite-transcendence-degree over
\(\mathbf Q\). But (4) tends to infinity. Algebraic independence over
\(K\) implies algebraic independence over \(\mathbf Q\), so
\(\operatorname{trdeg}_{\mathbf Q}\mathbf Q(T_0,T_1,\ldots)=\infty\), a
contradiction. \(\square\)

This strictly strengthens the safe consequences ``not D-finite'' and ``the
coefficient sequence is not P-recursive.'' It also implies that \(U\) belongs
to no finite-type ordinary Picard--Vessiot extension of \(\mathbf C(s)\),
because every element of such an extension is differentially algebraic over
the base differential field.

\subsubsection{Spectral formal-series corollary}

Let \(m_k\ge2\) be arbitrary and put

\[
 \sigma_k:=P_{m_k}(T_k^{-1}),
\]

where \(P_m\) is the universal inverse-trace polynomial of v2.3. For every
finite index set \(I\),

\[
 \operatorname{trdeg}_K K(\sigma_k:k\in I)
 =\operatorname{trdeg}_K K(T_k:k\in I),                                    \tag{GT-9}
\]

because \(\sigma_k\in K(T_k)\) and \(T_k\) is algebraic over
\(K(\sigma_k)\). Hence the formal series \(\sum_{k\ge0}\sigma_kz^k\) is
also differentially transcendental over \(\mathbf C(z)\). Convergence is
not needed for this statement: Mahler's theorem is formal.

\subsubsection{Exact overlap with Arreche (2013)}

C. E. Arreche, ``A Galois-theoretic proof of the differential transcendence
of the incomplete Gamma function,'' \emph{J. Algebra} \textbf{389} (2013), 119--127
(\href{https://arxiv.org/abs/1304.1917}{arXiv:1304.1917}), proves

\[
 \gamma(t,x)=\int_0^x u^{t-1}e^{-u}\,du
 \quad\text{is }\partial_t\text{-transcendental over }\mathbf C(t,x).       \tag{GT-10}
\]

Its setting is parameterized Picard--Vessiot theory: \(x\) is the main
differential variable, \(t\) is the parameter, and

\[
 \partial_x^2Y-\frac{t-1-x}{x}\partial_xY=0.                               \tag{GT-11}
\]

Arreche's criterion in fact applies to any nonconstant solution of (11), so
it also covers the bivariate upper incomplete Gamma after placing that
solution in the ambient differential extension.

The overlap and non-overlap are therefore precise:

\begin{itemize}
\tightlist
\item
  Both results assert differential transcendence in the \textbf{parameter/order}
  direction.
\item
  Arreche's theorem is unconditional and bivariate over \(\mathbf C(t,x)\),
  obtained from a parameterized Galois group.
\item
  Theorem 3 above first fixes the cut at \(x=1\) and then proves that the
  one-variable entire function \(s\mapsto\Gamma(s,1)\) is DT over
  \(\mathbf C(s)\), using the arithmetic independence theorem of {[}001{]} and
  Mahler's coefficient theorem.
\item
  The bivariate assertion (10) does \textbf{not formally specialize} to the
  fixed-cut assertion (7): differential transcendence need not survive
  specialization of a second variable. Conversely, (7) says nothing about
  the full PPV group in the \(x\)-direction.
\end{itemize}

Thus (7) is not a corollary of the cited Arreche theorem as stated. A final
priority claim would still require a dedicated search of the broader
shift-difference/hypertranscendence literature, because \(U\) also satisfies
\(U(s+1)=sU(s)+e^{-1}\).

\begin{center}\rule{0.5\linewidth}{0.5pt}\end{center}

\subsection{3. A half-density theorem for wholly transcendental tail divisors}

For \(k\ge0\), set

\[
 F_k(z):=\Ein(z)-T_k.
\]

\subsubsection{Lemma 4 (unique real tail root)}

There is a unique \(r_k\in(0,1)\) with

\[
 \Ein(r_k)=T_k,                                                             \tag{GT-12}
\]

and the sequence \((r_k)\) is strictly decreasing.

\paragraph{Proof}

With \(u=e^y\), put

\[
 I_k:=k!T_k=\int_0^\infty e^{-e^y}y^k\,dy.
\]

Integration by parts gives

\[
 (k+1)I_k
 =\int_0^\infty e^y e^{-e^y}y^{k+1}\,dy>I_{k+1},
\]

so \(T_{k+1}<T_k\). Also \(0<T_k\le T_0=E<\eta=\Ein(1)\).
The real function \(\Ein\) is strictly increasing because
\((1-e^{-x})/x>0\) for every real \(x\ne0\). This proves existence,
uniqueness, and strict decrease. \(\square\)

\subsubsection{Lemma 5 (only possible algebraic zero)}

The only possible algebraic zero of \(F_k\) is \(r_k\).

\paragraph{Proof}

There is only one real zero by strict monotonicity. If a nonreal algebraic
zero \(z\) existed, then \(\bar z\ne z\) would also be algebraic and

\[
 \Ein(z)=T_k=\Ein(\bar z).
\]

This is a polynomial relation between two distinct coordinates of the
jointly algebraically independent family in {[}001, Cor. 8.4{]}, a contradiction.
\(\square\)

\subsubsection{Theorem 6 (half-density all-zero theorem)}

Let

\[
 A_N:=\#\{0\le k\le N:r_k\in\overline{\mathbf Q}\}.
\]

Then for every \(N\ge1\),

\[
 \boxed{A_N\le2+\left\lfloor\frac N2\right\rfloor.}                         \tag{GT-13}
\]

Consequently at least

\[
 \boxed{\left\lceil\frac N2\right\rceil-1}                                 \tag{GT-14}
\]

of the divisors \(F_0,\ldots,F_N\) have \textbf{all} their zeros transcendental.
In density notation,

\[
 \liminf_{N\to\infty}
 \frac{\#\{0\le k\le N:\text{every zero of }F_k\text{ is transcendental}\}}
 {N+1}\ge\frac12.                                                           \tag{GT-15}
\]

\paragraph{Proof}

Let \(I\subset\{0,\ldots,N\}\) be the set of indices with algebraic
\(r_k\). The \(2|I|\) resonant coordinates

\[
 \bigl\{\Ein_{k+1}(1),\ \Ein_1(r_k):k\in I\bigr\}
\]

are pairwise distinct and jointly algebraically independent over \(K\) by
{[}001, Cor. 8.4{]}. Via (1) and \(g_k=c_k+T_k\), the invertible blockwise
linear transformation

\[
 (c_k,T_k)\longleftrightarrow(c_k,g_k)
\]

shows that the selected \(g_k\)'s are algebraically independent over
\(K\). Hence

\[
 |I|\le\operatorname{trdeg}_K K(g_0,\ldots,g_N)
 \le2+\left\lfloor\frac N2\right\rfloor,
\]

which proves (13). Lemma 5 then gives (14) and (15). \(\square\)

\subsubsection{A sharp level-zero fork}

The unique real zero at level zero is

\[
 r_0=\Ein^{-1}(E_1(1))\approx0.2321964241289549.
\]

Exactly one of the following holds:

\begin{enumerate}
\def\labelenumi{\arabic{enumi}.}
\tightlist
\item
  every zero of \(\Ein(z)-E_1(1)\) is transcendental;
\item
  \(r_0\) is algebraic, and \(\gamma\) and \(\delta\) are algebraically
  independent over \(K\).
\end{enumerate}

Indeed, in the second case \(\Ein(1)=\eta\) and \(\Ein(r_0)=E\) are jointly
algebraically independent. The invertible linear/scalar change

\[
 (\eta,E)\longleftrightarrow(\gamma=\eta-E,\delta=eE)
\]

proves the assertion. This is a genuine zero criterion, but not an
unconditional irrationality result: the arithmetic nature of \(r_0\) is
unknown.

\subsubsection{An Euler--Gompertz pencil strengthening}

For real algebraic \(t\in[0,1]\), put

\[
 c_t:=(1-t)\gamma+tE,
 \qquad \Ein(q_t)=c_t,\quad q_t\in(0,1).                                    \tag{GT-16}
\]

The elementary bounds \(E<e^{-1}<1/2<\gamma\) show that the \(c_t\)'s,
and hence the \(q_t\)'s, are distinct as \(t\) varies. Each divisor
\(\Ein-c_t\) has at most one algebraic zero, namely \(q_t\), and among the
entire algebraic pencil there is at most one algebraic \(q_t\). For if
\(q_s,q_t\) were algebraic with \(s\ne t\), the three distinct resonant
values \(\eta,c_s,c_t\) would be jointly algebraically independent, whereas

\[
 (t-s)\eta+(1-2t)c_s-(1-2s)c_t=0.                                          \tag{GT-17}
\]

Moreover \(q_{1/2}\) is unconditionally transcendental because
\(2\Ein(q_{1/2})-\Ein(1)=0\). If the one possible exceptional \(q_t\) is
algebraic, then \(t\ne1/2\), and the matrix taking
\((\gamma,E)\) to \((\eta,c_t)\) has determinant \(2t-1\ne0\). Therefore
that exception would again imply algebraic independence of \(\gamma\) and
\(\delta\) over \(K\).

\begin{center}\rule{0.5\linewidth}{0.5pt}\end{center}

\subsection{4. What the inverse traces do and do not add to the Euler motive}

Let \(P_m\in\mathbf Q[X]\) be the monic universal trace polynomial of v2.3,

\[
 \mathcal T_m(c)=P_m(c^{-1}),\qquad m\ge2,
\]

and recall the exact recovery identity

\[
 c^{-1}=144\mathcal T_4(c)-144\mathcal T_2(c)^2-14\mathcal T_2(c).          \tag{GT-18}
\]

\subsubsection{Exact spectral reformulations}

For distinct algebraic \(s,t\in[0,1]\) and arbitrary \(m,n\ge2\),

\[
 \boxed{
 \operatorname{trdeg}_K
 K\bigl(\mathcal T_m(c_s),\mathcal T_n(c_t)\bigr)
 =\operatorname{trdeg}_K K(\gamma,\delta)\in\{1,2\}.}                     \tag{GT-19}
\]

Each level is algebraic over any one of its nonconstant trace coordinates,
and the two distinct linear pencil levels recover \(\gamma,E\). Thus
\(\gamma,\delta\) are algebraically independent over \(K\) if and only if
any such pair of distinct pencil traces is algebraically independent. This
is an exact algebraic-matroid/spectral criterion, not a proof of either side.

For the Euler exponential motive with period matrix

\[
 P_\gamma=\begin{pmatrix}1&\gamma\\0&2\pi i\end{pmatrix},
\]

full comparison injectivity of

\[
 \overline{\mathbf Q}[a,a^{-1},b]\longrightarrow\mathbf C,
 \qquad a\mapsto2\pi i,\quad b\mapsto\gamma,                              \tag{GT-20}
\]

is equivalent to

\[
 \operatorname{trdeg}_{\overline{\mathbf Q}}
 \overline{\mathbf Q}(2\pi i,\gamma)=2.                                   \tag{GT-21}
\]

If \(S_m=P_m(\gamma^{-1})\), (18) gives

\[
 \boxed{
 (20)\text{ injective}
 \iff
 \operatorname{trdeg}_{\overline{\mathbf Q}}
 \overline{\mathbf Q}(2\pi i,S_2,S_4)=2.}                                 \tag{GT-22}
\]

Again, (22) is a birational re-encoding of the same unknown period
statement.

\subsubsection{The exact barriers}

\begin{enumerate}
\def\labelenumi{\arabic{enumi}.}
\item
  \textbf{Global DT does not isolate the constant term.} Equation (7) only says
  that infinitely much arithmetic complexity occurs throughout the tail.
  It is compatible with \(T_0=\delta/e\) being algebraic. Likewise (4)
  can be carried entirely by higher \(T_k\)'s.
\item
  \textbf{The gamma-jet upper field has enough unknown capacity.} The odd
  zeta values and \(\pi^2\) in (2) can absorb roughly half of the finite
  transcendence rank. No argument above forces \(\gamma\) to contribute
  a transcendence degree.
\item
  \textbf{Inverse traces are not automatically motivic periods.} The period
  algebra of a Tannakian category is not generally closed under inversion
  of an arbitrary numerical period coordinate. Formula (18) analytically
  recovers \(\gamma\), but it does not construct \(S_2,S_4\) as matrix
  coefficients of a tensor construction on \(M_\gamma\). Consequently it
  supplies no new comparison injectivity by itself.
\item
  \textbf{Degree-one fullness is already the hard \(\gamma\) problem.} The
  Euler motive is known to be a non-split extension
  \(0\to\mathbf Q(0)\to M_\gamma\to\mathbf Q(-1)\to0\). In this fixed
  rank-two setting, the assertion ``every numerical relation
  \(\lambda\gamma+\beta=0\) is motivic and forces splitting'' is logically
  equivalent to transcendence of \(\gamma\): non-splitting makes fullness
  imply transcendence, while transcendence leaves no such relation to
  lift.
\item
  \textbf{Bare Frobenius does not retain the extension.} The unfiltered
  rank-two Frobenius object splits automatically in the regime proved in
  {[}001{]}. Any successful arithmetic comparison must retain filtered,
  syntomic, or coherent adelic norm data and prove a numerical-to-motivic
  lifting theorem. The new inverse-root or Hankel identities do not supply
  that missing functorial comparison.
\item
  \textbf{Function-level differential Galois rigidity is a different axis.}
  The conclusion that \(U(s)\) lies in no finite Picard--Vessiot extension
  concerns variation in the parameter \(s\). It does not turn a numerical
  relation at \(s=0\) into a morphism of the Euler motive and therefore does
  not imply irrationality of \(\delta\), transcendence of \(\gamma\), or
  period-map injectivity.
\end{enumerate}

\subsubsection{One genuine motivic-size consequence}

No single fixed finite-dimensional Tannakian object with finite-type motivic
Galois group can contain the entire numerical tail tower \((T_k)_{k\ge0}\)
in the field generated by its period-matrix entries: that field has finite
transcendence degree, whereas (5) has infinite transcendence degree. Any
motivic realization of all tail jets must therefore use an increasing tower
or an infinite/pro-object. This is a structural consequence, not a new
statement about \(M_\gamma\)'s rank-two comparison map.

\begin{center}\rule{0.5\linewidth}{0.5pt}\end{center}

\subsection{5. Consolidated classification}

\small
\begin{longtable}[]{@{}
  >{\raggedright\arraybackslash}p{(\columnwidth - 4\tabcolsep) * \real{0.2727}}
  >{\raggedleft\arraybackslash}p{(\columnwidth - 4\tabcolsep) * \real{0.3636}}
  >{\raggedleft\arraybackslash}p{(\columnwidth - 4\tabcolsep) * \real{0.3636}}@{}}
\toprule\noalign{}
\begin{minipage}[b]{\linewidth}\raggedright
Claim
\end{minipage} & \begin{minipage}[b]{\linewidth}\raggedleft
Status from {[}001{]} + v2.3
\end{minipage} & \begin{minipage}[b]{\linewidth}\raggedleft
Arithmetic force on \(\gamma,\delta\)
\end{minipage} \\
\midrule\noalign{}
\endhead
\bottomrule\noalign{}
\endlastfoot
\(\operatorname{trdeg}_K K(T_0,\ldots,T_N)\ge\lceil N/2\rceil-1\) & proved & none individually \\
\(\Gamma(s,1)\) DT over \(\mathbf C(s)\) & proved via Mahler & none individually \\
no finite ordinary PV extension contains \(\Gamma(s,1)\) & proved & none individually \\
lower density \(\ge1/2\) of tail divisors with all zeros transcendental & proved & none individually \\
algebraic \(r_0\Rightarrow \gamma,\delta\) jointly AI over \(K\) & proved conditional fork & very strong if trigger is settled \\
at most one algebraic root in the Euler--Gompertz pencil & proved & an exception forces joint AI \\
two trace coordinates recover an Euler/pencil level & proved & exact criterion only \\
full Euler period-map injectivity & open & equivalent to \(\pi,\gamma\) algebraic independence \\
\(\gamma\) transcendental from motivic non-splitting alone & not proved & needs numerical-to-motivic fullness \\
\(\delta\) irrational/transcendental & not proved & tail rank does not select \(T_0\) \\
\end{longtable}
\normalsize

The strongest transfer package is therefore the combination of (i) a
\textbf{fixed-cut differential-transcendence theorem}, (ii) a
\textbf{half-density all-zero theorem}, and (iii) sharp \textbf{conditional zero/trace
criteria}. None closes Euler's or Gompertz's individual arithmetic problem
without one further trigger: an algebraic exceptional inverse root, or a
numerical-to-motivic comparison theorem.
